\chapter{Увод}

Развој метода машинског и дубоког учења омогућио је анализу сложених појава чија
структура измиче традиционалним статистичким приступима. У области финансија,
међутим, примена ових метода наилази на границу која није алгоритамске, него
структурне природе. Гу, Кели и Сју \parencite{guKellyXiu2020} упоредили су читав
репертоар метода машинског учења на скупу од преко деветсто предиктора и утврдили да
и најбоље међу њима објашњавају мање од пола процента варијансе месечних приноса.
Тај налаз указује на то да правац кретања цена није проблем са скривеним решењем,
већ проблем у коме сигнала готово да нема.

Постоји, међутим, изразита асиметрија у предвидљивости финансијских временских
серија. Док је аутокорелација приноса практично једнака нули, аутокорелација
квадрираних приноса снажно је позитивна и споро опада, остајући значајна и након више
десетина трговачких дана. Другим речима, не зна се да ли ће цена расти или падати,
али се са знатном тачношћу може предвидети колико ће се снажно кретати. Ова асиметрија
има економско објашњење: предвидљивост правца сама себе уништава, јер је арбитражери
искоришћавају чим се појави, док предвидљивост интензитета не носи директну опкладу
којом би била избрисана. Реализована волатилност, мерена сумом квадрата интрадневних
логаритамских приноса, стога представља ретку величину у финансијама која има стваран
и постојан сигнал, а истовремено и величину на којој почивају регулаторно обавезни
поступци управљања ризиком, попут израчунавања вредности при ризику и вредновања
опција.

Модели за предвиђање реализоване волатилности развијали су се у две паралелне линије.
Економетријску линију утемељио је Корси \parencite{corsi2009} хетерогеним
ауторегресивним моделом, који дуготрајно памћење волатилности апроксимира каскадом
компоненти на дневном, недељном и месечном хоризонту, инспирисаном хипотезом о
хетерогеном тржишту, дакле идејом да учесници реагују различитим брзинама. Проширења
овог модела уводе разлагање на континуирану и скоковиту компоненту, док Патон и Шепард
\parencite{pattonSheppard2015} показују да је негативна семиваријанса знатно значајнија
за будућу волатилност од позитивне, и да негативни скокови воде вишој, а позитивни
нижој будућој волатилности. Друга линија примењује методе машинског учења: Кристенсен,
Сигор и Велијев \parencite{christensen2023} налазе да алгоритми машинског учења
надмашују HAR породицу чак и када су им једини предиктори исте дневне, недељне и
месечне доцње, уз највеће добитке на дужим хоризонтима.

Актуелна расправа концентрисана је на трансформерске архитектуре и у њој постоји
отворена противречност. Студије на америчким индексима налазе да трансформерски модели,
укључујући PatchTST, доследно постижу најниже грешке предвиђања, док радови усмерени на
окружења вођена скоковима и тешким реповима расподеле тврде да перформансе произлазе из
архитектурне индуктивне пристрасности, а не из капацитета модела, те да једноставни
линеарни модели у таквим условима надмашују трансформере. Свеобухватан преглед области
\parencite{leushuisPetkov2026}, који обухвата поређење модела у литератури од 2000. до
прве половине 2024. године, потврђује да питање остаје неразрешено.

Оно што целу ову расправу повезује јесте заједничка прећутна претпоставка: реализоване
мере третирају се као фиксни, унапред дати улази, а варира се искључиво архитектура
модела. Та претпоставка је проблематична, јер реализована волатилност није посматрана,
него изведена величина. Свака њена вредност резултат је низа методолошких одлука: избора
фреквенције узорковања на којој се рачунају приноси, избора естиматора отпорног на
скокове којим се процењује континуирани део волатилности, и избора статистичког теста
којим се дан проглашава даном са скоком. Ниједна од тих одлука није неутрална: превисока
фреквенција узорковања уводи микроструктурни шум који систематски надувава меру, а
различити тестови за детекцију скокова дају битно различите скупове детектованих
догађаја. Да грешка мерења заиста утиче на прогнозу, показали су Болерслев, Патон и
Кведвлиг \parencite{bpq2016} HARQ моделом, који коефицијент уз дневни регресор интерагује
са реализованом квартичношћу управо да би узео у обзир временски променљиву грешку
мерења. Тај модел, међутим, грешку мерења моделује унутар једног модела, уместо да је
третира као димензију која се може експериментално варирати.

Недостаје, дакле, систематско мерење изолованог доприноса самог начина мерења, и то не
само на једном моделу, него кроз цео распон архитектура. Овај рад се бави тим питањем
на панелу америчких акција и кроз потпун факторски експеримент.

\section{Циљ рада}

Општи циљ рада јесте да се измери колики удео разлика у тачности предвиђања реализоване
волатилности потиче од избора модела, а колики од одлука донетих при мерењу улазних
величина.

Конкретни циљеви рада су:

\begin{itemize}
\item имплементирати сопствени систем за конструкцију реализованих мера, у коме се
      дизајн мерења може мењати по димензијама, без измене осталих делова поступка;
\item спровести факторски експеримент у коме се тридесет шест дизајна мерења укршта са
      седамнаест модела, кроз три хоризонта предвиђања и једанаест валидационих прозора;
\item измерити тачност предвиђања функцијом губитка робусном на шум у проксију, уз
      тестове статистичке значајности и скуп поверења модела;
\item бројчано раздвојити удео варијансе губитка који се приписује дизајну мерења од
      удела који се приписује избору модела, помоћу декомпозиције варијансе;
\item утврдити да ли се рангирање модела мења при преласку са статистичке на економску
      евалуацију, путем тестова вредности при ризику;
\item утврдити да ли побољшања приписана информацији о скоковима опстају када се
      индикатори скокова насумично пермутују, чиме се ефекат информације разликује од
      ефекта структуре модела.
\end{itemize}

\section{Методологија}

За потребе рада развијен је сопствени систем за конструкцију реализованих мера и
оцењивање модела, у програмском језику Python. Интрадневне свећице на једноминутној
резолуцији преузимају се из берзанског извора и чисте према стварном трговачком
календару, након чега се за сваку комбинацију дизајна мерења рачунају реализоване мере
и спроводи условна декомпозиција на континуирану и скоковиту компоненту.

Дизајн мерења вари́ра се кроз три независне димензије: фреквенцију узорковања, естиматор
отпоран на скокове и тест за детекцију скокова, чијим се укрштањем добија тридесет шест
дизајна. Сваки од њих представља један легитиман, у литератури заступљен начин
конструисања истих улазних величина. Скуп модела обухвата наивну основу, HAR породицу и
њена проширења, дугомеморијске и мултипликативне економетријске моделе, моделе засноване
на стаблима одлучивања, трансформерску архитектуру и темељне моделе временских серија
примењене без обучавања.

Оцењивање се спроводи кроз усидрену унапредну валидацију, којом се искључује свако
цурење информација из будућности. Тачност се мери QLIKE функцијом губитка, изабраном
зато што је Патон \parencite{patton2011} показао да већина уобичајених функција губитка
даје погрешно рангирање модела када је прокси за стварну волатилност зашумљен, а да је
QLIKE међу малобројним функцијама робусним на ту појаву. Значајност разлика утврђује се
упареним тестом предиктивне тачности и скупом поверења модела, уз контролу стопе лажних
открића, а доприноси појединачних фактора раздвајају се декомпозицијом варијансе.

Мере, правила искључивања и план статистичке анализе записани су пре покретања
експеримента. Сви параметри који нису предмет варирања фиксирани су унапред и наведени
уз образложење. Детаљи су описани у поглављу~\ref{ch:metodologija}.

\section{Хипотезе}

Полазна теза рада јесте да дизајн мерења представља експерименталну димензију
равноправну избору модела, те да део објављених разлика међу моделима може потицати од
начина припреме података. Да би се та теза могла проверити, рашчлањена је на пет
хипотеза:

\begin{itemize}
\item \textbf{Х1:} Додавање компоненти скокова као улазних променљивих побољшава
      тачност предвиђања у односу на моделе који их не користе.
\item \textbf{Х2:} То побољшање израженије је на кратком хоризонту, где скокови чине
      већи део варијације, него на дугом, где доминира перзистентност волатилности.
\item \textbf{Х3:} Варијација у дизајну мерења објашњава удео варијансе губитка упоредив
      са уделом који објашњава избор модела.
\item \textbf{Х4:} Рангирање модела мења се при преласку са статистичке мере тачности на
      економску евалуацију путем вредности при ризику.
\item \textbf{Х5:} Побољшања приписана скоковима нестају када се индикатори скокова
      насумично пермутују кроз време.
\end{itemize}

Централна хипотеза јесте Х3, јер се непосредно односи на постављени циљ рада и
проверава се декомпозицијом варијансе дневних губитака. Хипотезе Х1 и Х2 проверавају да
ли информација о скоковима уопште носи сигнал, без чега налаз о Х3 не би био
интерпретабилан. Хипотеза Х4 испитује да ли закључак опстаје при промени мерила, а
хипотеза Х5 представља плацебо проверу којом се утврђени ефекат разликује од артефакта
структуре модела.

Ниједан исход није неуспех истраживања. Ако се Х3 потврди, следи да дизајн мерења
заслужује да буде наведен и образложен у свакој студији о предвиђању волатилности,
једнако као што се наводи архитектура модела. Ако се не потврди, рад пружа емпиријски
основ за супротну тврдњу, дакле да су реализоване мере довољно робусне да избор међу
њима не утиче битно на поређење модела. Та тврдња се тренутно претпоставља, али није
проверена. Разлике у ефекту између појединих режима и хоризоната анализирају се
експлоративно, без унапред формулисане хипотезе.

\section{Структура рада}

Рад је организован у шест поглавља. У другом поглављу изложене су теоријске основе:
конструкција реализованих мера и њихово разлагање на континуирану и скоковиту
компоненту, два експериментална фактора која рад укршта, изведене променљиве које чине
улаз модела, као и валидациони протокол и апарат за евалуацију. Треће поглавље описује
методологију: извор и припрему података, дизајн експеримента са фиксираним параметрима,
мере, правила искључивања и план статистичке анализе. У четвртом поглављу приказани су
поставка и резултати експеримента, њихова дискусија у односу на хипотезе и ограничења
рада. Пето поглавље сумира главне налазе и наводи правце даљег истраживања.
